% Pour toute question sur l'utilisation de cette classe, écrivez à jerome.soucy(at)ulaval.ca
% Nécessite la classe Evaluation.sty dans le même dossier ou dans le PATH de votre système
% Nécessite le fichier logo_ul.pdf dans le même dossier un dans le PATH de votre système
\def\Ponderation{8}
\def\SigleNumero{STT-2100}
\def\NomCours{Régression}
%\def\DateEvaluation{2 mars 2023}
\def\DateRemise{1er octobre 2026}
\def\HeureRemise{23h55}
\def\HeureEvaluation{}
\def\Session{Été 2021}
\def\NomEvaluation{Devoir 1}

% Ne rien mettre entre les accolades si l'enseignant (ou l'enseignante) est aussi le coordonnateur (ou la coordonnatrice)
\def\Coordonnatrice{}
\def\Coordonnateur{}

% Ne rien mettre entre les accolades s'il y a plus d'une section
\def\Enseignant{}
\def\Enseignante{Audrey-Anne Vallée}

% Ne rien mettre entre les accolades s'il s'agit d'un cours à section unique
% Mettre le nom de chacun des enseignants et enseignantes entre les accolades
% Une case à cocher sera générée pour chacune des sections
\def\SectionA{}
\def\SectionB{}
\def\SectionC{}
\def\SectionD{}


\def\Directives{
\begin{enumerate}
\item Le devoir compte \numquestions ~questions pour un total de \numpoints{}~points.
\item À faire individuellement ou en équipe de deux personnes. 

\item Vous devez remettre un seul fichier PDF contenant vos réponses aux questions 1 à \numquestions. Le PDF avec vos solutions doit être produit à partir d’un logiciel comme Word, LaTeX, etc. La numérisation de solutions manuscrites ne sera pas acceptée. Le fichier est à remettre sur Gradescope, à l'aide des instructions données au point 4. sur cette page.

\item Remises sur \textit{Gradescope}:
\begin{itemize}
\item Rendez-vous au www.gradescope.com;
\item Connectez-vous en utilisant votre adresse courriel de l'Université (idul@ulaval.ca) ainsi que votre mot de passe habituel. Il faut faire «Log in», utiliser «School credentials» et chercher Université Laval dans la liste . Vous devriez déjà avoir un compte, par le biais de l’université. 
\item Une fois sur le site, choisissez STT-2100;
\item En sélectionnant le Devoir 1, vous ajoutez le PDF qui contient vos réponses;
\item Vous devez identifier les différentes sous-questions sur les pages de votre devoir. Si cette étape n'est pas bien faite, une pénalité de 5 points peut être attribuée.
\end{itemize}

\item Il est \textit{\textbf{nécessaire}} de justifier toutes les démarches. Une réponse qui n'a pas une démarche \textit{\textbf{complète}} ne pourra avoir tous les points. Vous ne pouvez pas utiliser les résultats du cours directement, à moins que ce soit indiqué.
\item La qualité de la langue sera évaluée, incluant l'orthographe, la grammaire et la syntaxe. Un maximum de 10 points pourra être retranché pour mauvaise qualité du français. 
\item Les travaux remis en retard recevront la note 0.
%\item Veuillez relire le point 7, il n'y aura aucune exception.
\end{enumerate}
}

% Veuillez indiquer le type d'évaluation entre crochets
% Examen : option à choisir pour un examen présentiel qui sera corrigé de manière classique. Un tableau des points sera affiché sur la page couverture.
% Gradescope : option à choisir pour un examen qui sera corrigé via Gradescope. Aucun tableau de pointage n'apparaîtra et il y a quelques différences au niveau de la mise en page
% TPS :
% TPF :
% Devoir :
\documentclass[Devoir]{Evaluation}
\usepackage{multicol,graphicx}
\begin{document}

\fbox{%
    \parbox{\textwidth}{%

Pour certains exercices, nous référerons au modèle de régression linéaire simple donné par

\begin{equation}\label{mod:reglin}
    Y_i = \beta_0\,+\,\beta_1\, x_i \,+\, \varepsilon_i\qquad \text{ pour } i=1,\ldots,n.
\end{equation}

Dans ce modèle, les 4 postulats sont respectés et nous avons $\varepsilon_i \overset{iid}{\sim} \mathcal{N}(0, \sigma^2), \text{ pour } i=1,\ldots,n$.
    }%
}




\begin{questions}  

\vspace{1cm}
\question

Considérons le modèle de régression donné par l'équation \eqref{mod:reglin} au début du document. Dans cet exercice, on s'intéresse à l'estimation de $\sigma^2$. \\

On a que $\varepsilon_i \overset{iid}{\sim} \mathcal{N}(0, \sigma^2), \text{ pour } i=1,\ldots,n$. 
On pourrait croire que les résidus suivent également une loi normale de moyenne $0$ et de variance $\sigma^2$ et sont indépendants. Ce n'est toutefois pas le cas. Montrons que
\begin{equation*}
\widehat{\varepsilon}_j \sim \mathcal{N} \left\{0, \, (1-h_{jj})\sigma^2\right\}, \quad \mbox{ où } h_{jj}=\frac{1}{n}+\frac{(x_j - \bar x)^2}{S_{xx}}.
\end{equation*}
Exprimons les résidus sous la forme d'une combinaison linéaire des $Y_i$,
\begin{eqnarray*}
 \widehat{\varepsilon}_j &=& Y_j - \widehat{Y}_j\\
 &=& Y_j - \widehat{\beta}_0 - \widehat{\beta}_1 x_j\\
 &=& Y_j - \sum_{i=1}^n \left(\frac{1}{n} - a_i \overline{x}\right) Y_i - \left(\sum_{i=1}^n a_i Y_i\right) x_j \\
 &=& \sum_{i=1}^n d_i Y_i,
\end{eqnarray*}
où
\begin{eqnarray*}
a_i &= & \frac{\displaystyle x_i-\bar x}{\displaystyle S_{xx}}, \quad
d_i  =  \left\{
 \begin{array}{lc}
 \displaystyle 1-\frac{1}{n} + a_j (\overline{x}-x_j) & \mbox{ si } i=j, \\[4mm]
  \displaystyle     -\frac{1}{n} + a_i (\overline{x}-x_j) & \mbox{ si } i\neq j.
     \end{array}
     \right.
\end{eqnarray*}
Ainsi, les résidus suivent une distribution normale. Notons que
\begin{eqnarray*}
\sum_{i=1}^n d_i =   0, \quad \sum_{i=1}^n d_i x_i =   0, \quad \text{et que} \quad \sum_{i=1}^n d_i^2 = 1-\left\{\frac{1}{n}+\frac{(\bar x-x_j)^2}{S_{xx}}\right\}.
\end{eqnarray*}

\begin{parts}
    \part[5] Montrez que $\mathrm{E}\left(\widehat{\varepsilon}_j\right) = 0.$
    \part[5] Montrez que $\mathrm{Var}\left(\widehat{\varepsilon}_j\right) =  (1-h_{jj})\sigma^2.$
    \part[10] En utilisant le résultat la distribution des résidus, montrer que $ \displaystyle \widehat\sigma_{MV}^2=\frac{\displaystyle 1}{\displaystyle n}\sum\limits_{i=1}^n \widehat\varepsilon_i^2$
est un estimateur  biaisé de  $\sigma^2$. 
\part[5] En utilisant le résultat de (c), montrez que $ \displaystyle \widehat\sigma^2=\frac{\displaystyle 1}{\displaystyle n-2}\sum\limits_{i=1}^n \widehat\varepsilon_i^2$ est un estimateur sans biais de $\sigma^2$.
\end{parts}


\begin{solution}
    \begin{parts}
        \part \begin{align*}
        \mathrm{E}\left(\widehat{\varepsilon}_j\right) 
        &= \sum_{i=1}^n d_i \mathrm{E}\left(Y_i\right) \cr
        &= \sum_{i=1}^n d_i (\beta_0 + \beta_1 x_i) \cr
        &= \beta_0\sum_{i=1}^n d_i  + \beta_1 \sum_{i=1}^n d_i x_i \cr
        &=0.
        \end{align*}
        \part 
        \begin{align*}
        \mathrm{Var}\left(\widehat{\varepsilon}_j\right) 
        &= \sum_{i=1}^n d_i^2 \mathrm{Var}\left(Y_i\right)\\
        &= \left[1-\left\{\frac{1}{n}+\frac{(\overline{x}-x_j)^2}{S_{xx}}\right\}\right]\sigma^2 \\
        &=(1-h_{jj})\sigma^2.
        \end{align*}
        \part  \begin{align*}
            \mathrm{E}\left(\widehat{\sigma}_{MV}^2\right)
            &=\mathrm{E}\left(\frac{1}{n} \sum\limits_{i=1}^n \widehat\varepsilon_i^2\right) \cr
            &= \frac{1}{n} \sum\limits_{i=1}^n \mathrm{E}\left(\widehat\varepsilon_i^2\right) \cr
             &= \frac{1}{n} \sum\limits_{i=1}^n \left[ \mathrm{Var}\left(\widehat\varepsilon_i\right) + \mathrm{E}\left(\widehat\varepsilon_i\right)^2 \right]\cr
             &= \frac{1}{n} \sum\limits_{i=1}^n \left[ \left(1- \frac{1}{n} - \frac{(x_i - \overline{x})^2}{S_{xx}}\right)\sigma^2+ 0\right]\cr
             &= \frac{\sigma^2}{n} \sum\limits_{i=1}^n \left[ 1- \frac{1}{n} - \frac{(x_i - \overline{x})^2}{S_{xx}}\right]\cr
             &= \frac{\sigma^2}{n}  \left[ n- 1 - 1\right]\cr
             &= \sigma^2 \frac{n-2}{n}
        \end{align*}
        \part \begin{align*}
            \mathrm{E}\left(\widehat{\sigma}^2\right)
            &= \mathrm{E}\left(\frac{n}{n-2}\widehat{\sigma}_{MV}^2\right) 
            = \frac{n}{n-2} \mathrm{E}\left(\widehat{\sigma}_{MV}^2\right) 
            = \frac{n}{n-2}\frac{n-2}{n}\sigma^2 
                = \sigma^2.
            \end{align*}
    \end{parts}
\end{solution}







\question

Soit le modèle de régression donné par l'équation \eqref{mod:reglin} au début du document. Nous tentons d'estimer $\beta_0$ et~$\beta_1$ dans ce modèle à l'aide de $\widehat{\beta}_0^*$ et $\widehat{\beta}_1^*$ que nous définissons comme les valeurs qui minimisent la fonction:

\begin{equation*}
   Q(\beta_0, \beta_1)= \sum_{i=1}^n(Y_i-\beta_0- \beta_1x_i)^2 + \lambda \beta_1 ^2,
\end{equation*}

où $\lambda \in \mathbb{R}_+$. (Aux sous-questions (b) et (d), vous pouvez utiliser directement  les résultats du chapitre 3.)
\begin{parts}
    \part[10] Donnez les expressions de $\widehat{\beta}_0^*$ et de $\widehat{\beta}_1^*$.
    \part[4] Calculez le biais de $\widehat{\beta}_1^*$.
    \part[1] Pour quelle(s) valeur(s) de $\lambda$ a-t-on que $\widehat{\beta}_1^*$ est sans biais pour $\beta_1$?
    \part[8] Pour quelle(s) valeur(s) de $\lambda$ a-t-on $\text{Var}\left(\widehat{\beta}_1^*\right)<\text{Var}\left(\widehat{\beta}_1\right)$, où $\widehat{\beta}_1$ est l'estimateur des moindres carrés ordinaires?
    \part[2] Quel serait l'intérêt d'utiliser une valeur de $\lambda$ qui induit un biais dans notre estimation de~$\beta_1$?
    
\end{parts}

\begin{solution}
    \begin{parts}
        \part \begin{align*}
            \frac{\partial Q}{\partial \beta_0}
            &= -2 \sum (Y_i-\beta_0- \beta_1x_i) \cr
            %
            \frac{\partial Q}{\partial \beta_1}
            &= -2 \sum (Y_i-\beta_0- \beta_1x_i)x_i + 2\lambda \beta_1 
        \end{align} 
        \begin{align*}
            \sum (Y_i-\widehat{\beta}_0^* - \widehat{\beta}_1^* x_i) 
            &= 0 \cr
            n\overline{Y} - n\widehat{\beta}_0^* - n\widehat{\beta}_1^* \overline{x}
            &= 0 \cr
            \overline{Y} - \widehat{\beta}_1^* \overline{x}
            &=\widehat{\beta}_0^* \cr
            %
            \sum (Y_i-\widehat{\beta}_0^* - \widehat{\beta}_1^* x_i) x_i -\lambda \widehat{\beta}_1^* 
            &= 0 \cr
            \sum Y_ix_i - n\widehat{\beta}_0^*\overline{x} - \widehat{\beta}_1^* \sum x_i^2 -\lambda \widehat{\beta}_1^*
            &= 0 \cr
            \sum Y_ix_i - n(\overline{Y} - \widehat{\beta}_1^* \overline{x})\overline{x} - \widehat{\beta}_1^* \sum x_i^2 -\lambda \widehat{\beta}_1^* 
            &= 0 \cr
            \sum Y_ix_i - n\overline{Y}\overline{x} + n\widehat{\beta}_1^* \overline{x}^2 - \widehat{\beta}_1^* \sum x_i^2 -\lambda \widehat{\beta}_1^* 
            &= 0 \cr
            S_{xy}
            &=  \widehat{\beta}_1^* S_{xx} +\lambda \widehat{\beta}_1^* \cr
            \frac{S_{xy}}{S_{xx} + \lambda}
            &=  \widehat{\beta}_1^* *
            \cr
        \end{align} 
        \part  $$ \mathrm{E} (S_{xy}/S_{xx}) = \beta_1  \rightarrow \mathrm{E} (S_{xy}) = \beta_1 S_{xx}  $$
        $$ \mathrm{E} (S_{xy}/(S_{xx}+\lambda)) = \mathrm{E} (S_{xy})/(S_{xx}+\lambda) = \beta_1 S_{xx} / (S_{xx}+\lambda) $$
        \part $\lambda = 0$
        \part $$ \mathrm{Var} (S_{xy}/S_{xx}) = \sigma^2/S_{xx}  \rightarrow (1/S_{xx}^2)\mathrm{Var} (S_{xy}) = \sigma^2/S_{xx}  \rightarrow \mathrm{Var} (S_{xy}) = \sigma^2 S_{xx}  $$
         $$ \mathrm{Var} (S_{xy}/(S_{xx}+\lambda)) = \mathrm{V} (S_{xy})/(S_{xx}+\lambda)^2 = \sigma^2 S_{xx}  / (S_{xx}+\lambda)^2 $$
         $$  \frac{\sigma^2 S_{xx}}{(S_{xx}+\lambda)^2} \leq \frac{\sigma^2}{S_{xx}} \;\text{si}\;  
         \frac{ S_{xx}^2}{(S_{xx}+\lambda)^2} \leq 1 \;\text{si}\; 
         S_{xx}^2 \leq (S_{xx}+\lambda)^2 \;\text{si}\; 
         S_{xx}^2 \leq S_{xx}^2 + 2S_{xx}\lambda +\lambda^2 
         \;\text{si}\;  
         0 \leq  2S_{xx}\lambda +\lambda^2$$
         pour $\lambda > 0.$
         \part Diminuer l'EQM, avoir une plus petite variance.
    \end{parts}
\end{solution}




\vspace{1cm}
\question 
Soit le modèle de régression donné par l'équation \eqref{mod:reglin} au début du document. 
Nous supposons que les $x_i$ sont des variables binaires, ne pouvant prendre que les valeurs $0$ ou $1$.
Notons
$$  n_0=\sum_{i=1}^n (1-x_i); \quad
    n_1=\sum_{i=1}^n x_i; \quad
    \overline{Y}_0=\frac{1}{n_0}\sum_{i=1}^n (1-x_i)Y_i; \quad
    \overline{Y}_1=\frac{1}{n_1}\sum_{i=1}^n x_iY_i.$$


\begin{parts}
    \part[4] Dans vos mots, dites ce que représentent les quantités $n_0$, $n_1$, $\overline{Y}_0$ et $\overline{Y}_0$.
    \part[8] Montrez que l'estimateur des moindres carrés pour $\beta_1$ est donné par  $\widehat{\beta}_1=\overline{Y}_1-\overline{Y}_0 $. Vous pouvez utiliser le résultat $\widehat{\beta}_1= \displaystyle\frac{S_{xY}}{S_{xx}}$. (Indice: Commencez par écrire $\overline{Y}$ en fonction de $\overline{Y}_0$ et $\overline{Y}_1$, et par écrire $\overline{x}$ en fonction de $n_1$ et $n$).
    \part[3] Montrez que l'estimateur des moindres carrés pour $\beta_0$ est donné par  $\widehat{\beta}_0=\overline{Y}_0$. Vous pouvez utiliser le résultat $\widehat{\beta}_0= \overline{Y}-\widehat{\beta}_1\overline{x}$.
\end{parts}

\begin{solution}
    \begin{parts}
        \part \begin{itemize}
            \item $n_0$: le nombre d'individus ayant $x=0$.
            \item $n_1$: le nombre d'individus ayant $x=1$.
            \item $\overline{Y}_0$: la moyenne de $Y$ des individus ayant $x=0$.
            \item $\overline{Y}_1$: la moyenne de $Y$ des individus ayant $x=1$.
        \end{itemize}
        \part $ \overline{Y} = \frac{1}{n}\sum Y_i = \frac{1}{n} \left( \sum (1-x_i)Y_i + \sum x_iY_i \right) = \frac{1}{n}(n_0\overline{Y}_0 + n_1\overline{Y}_1)$ et $\overline{x}=n_1/n$
        \begin{align*}
            \widehat{\beta}_1
            &= \displaystyle\frac{S_{xY}}{S_{xx}}
            = \frac{\sum x_iy_i - n\overline{x}\overline{Y}}{\sum x_i^2 - n \overline{x}^2} \\
            &= \frac{n_1\overline{Y}_1 - \frac{n n_1}{n} \frac{1}{n}(n_0\overline{Y}_0 + n_1\overline{Y}_1)}{n_1 - n \frac{n_1^2}{n^2}}\\
            &= \frac{n_1\overline{Y}_1 - \frac{n_1}{n}(n_0\overline{Y}_0 + n_1\overline{Y}_1)}{n_1 -  \frac{n_1^2}{n}} \\
            &= \frac{n_1\overline{Y}_1 - \frac{n_1}{n}(n_0\overline{Y}_0 + n_1\overline{Y}_1)}{\frac{n_1}{n}(n- n_1)} \\
             &= \frac{n\overline{Y}_1 - (n_0\overline{Y}_0 + n_1\overline{Y}_1)}{(n- n_1)} \\
             &= \frac{(n-n_1)\overline{Y}_1 - n_0\overline{Y}_0 }{n_0} \\
             &= \overline{Y}_1 - \overline{Y}_0 
             \end{align*}
             \part 
        \begin{align*}
            \widehat{\beta}_0 &= \overline{Y}-\widehat{\beta}_1\overline{x} \cr
            &= \frac{1}{n}(n_0\overline{Y}_0 + n_1\overline{Y}_1) - (\overline{Y}_1 - \overline{Y}_0 ) \frac{n_1}{n} \\
            &= \frac{n_0}{n} \overline{Y}_0 + \frac{n_1}{n} \overline{Y}_1 - \frac{n_1}{n} \overline{Y}_1 +  \frac{n_1}{n} \overline{Y}_0 \\
            &= \overline{Y}_0
        \end{align*}
    \end{parts}
\end{solution}


\vspace{1cm}
\question\label{prog}

Vous devez choisir un jeu de données qui existe déjà dans \texttt{R}. Il peut s’agir d’un jeu de données inclus dans \texttt{R} de base ou fourni avec un package \texttt{R}. Vous ne devez donc pas créer votre propre jeu de données ni importer des données provenant d’une source externe.

Le jeu de données choisi doit contenir au moins deux variables quantitatives (non catégorielles) permettant de réaliser une régression linéaire simple dans un contexte pertinent. Au plus tard le jeudi 24 septembre 2026, vous devez m’envoyer un courriel dans lequel vous indiquez :
\begin{itemize}
    \item le nom du jeu de données;
    \item le nom du package dans lequel il se trouve, s’il y a lieu;
    \item les deux variables que vous souhaitez étudier.
\end{itemize}
Votre choix doit être approuvé avant que vous commenciez le travail. Deux équipes ne peuvent pas utiliser le même jeu de données. Les jeux de données approuvés et les équipes correspondantes seront inscrits dans un tableau sur le site du cours, sur la page du Devoir 1. Consultez ce tableau avant de soumettre votre choix.

Vous devez répondre aux questions suivantes,  en utilisant $\alpha = 0,01$ pour toutes vos inférences. Identifiez clairement la question associée à votre réponse (par exemple (i) a)). Donnez les réponses aux questions et les tableaux de résultats obtenus avec votre logiciel qui sont nécessaires pour
répondre aux questions. Les figures et tableaux inclus dans vos réponses doivent être présentés de façon professionnelle. Ils doivent avoir des titres clairs et pertinents. Chaque figure présentée doit être discutée dans vos réponses. 

\begin{parts}
	\part[9] {\bf Description du jeu de données et de l'objectif.}
\begin{enumerate}[(i)]
\item Identifiez et définissez les variables de votre modèle de RLS.%1
%\item Mentionner le nombre et la source des observations.
\item Décrivez {\it brièvement} la distribution de chacune des variables (graphiques et statistiques descriptives). Commentez la distribution conjointe des variables à l'aide d'un diagramme de dispersion. Votre choix de jeu de données et de variables doit être approprié dans un contexte de régression linéaire simple. (8 lignes maximum, sans compter les figures). %6
\item Décrire l'objectif de l'analyse de régression. (4 lignes maximum).%2
\end{enumerate}
	\part[16] {\bf Présentation des résultats.}
\begin{enumerate}[(i)]
	\item Donnez les estimations des param\`etres ${\beta}_0$, ${\beta}_1$, et ${\sigma}^2$. Écrivez l'expression de la droite de régression estimée.  %4
	\item Fournissez des intervalles de confiance pour les paramètres ${\beta}_0$ et ${\beta}_1$ et une interprétation en contexte. (8 lignes maximum pour l'interprétation). %5
	\item Discutez de la signification statistique des coefficients ${\beta}_0$ et ${\beta}_1$ à l'aide de tests $t$. (6 lignes maximum). %4
	\item Donnez la veleur du $R^2$. Interprétez cette valeur dans le contexte. Commentez la qualité du lien linéaire entre les variables. (6 lignes maximum pour la discussion). %3
\end{enumerate}
  \part[10] Copiez dans le PDF le code \texttt{R} utilisé pour répondre aux questions (a) et (b). Un programme qui ne fonctionne pas ou qui est mal annoté entraînera une pénalité. %
\end{parts}

\end{questions}


\end{document}




\vspace{0.3in}
\textbf{Question 4 (10 pts)} (Design d'expérience)\\

On s'intéresse à un médicament sensé réduire le nombre de crises chez les individus atteints d'epilepsie. Vous devez aider à planifier une expérience pour étudier l'effet de la dose du médicament sur le nombre de crises. Puisque vous vous intéressez à la relation entre deux variables quantitatives, vous planifiez analyser les données avec un modèle de régression linéaire simple. \\

Soit $x_1, \ldots, x_n$ les valeurs de la dose pour les 20 individus de votre étude.\\ Pour chacun, vous pouvez attribuer n'importe quelle dose entre 2 mg et 21 mg. \\

\begin{itemize}
%\item [(a)] Pour le modèle de régression linéaire simple $Y_i = \beta_0 + \beta_1 x_i + \varepsilon_i$, prouvez que la variance de $\beta_0$ est minimale si les $x_i$ sont tels que $\bar{x}$ = 0.
\item [(a)] (5 pts) Quelle est la meilleure fa\c{c}on de choisir les valeurs $x_1,\ldots,x_n$ pour minimiser la variance de votre estimateur $\widehat{\beta_1}$? Justifiez votre réponse.
\vspace{0.1 in}
\item [(b)] (5 pts) Pourquoi ne serait-ce pas vraiment une bonne idée de choisir les valeurs de $x_i$ de cette façon, même si cela minimise la variance de votre estimateur?
\end{itemize}

\end{document} 

\question 
Soit le modèle de régression donné par l'équation \eqref{mod:reglin} au début du document. Nous tentons d'estimer $\beta_0$ et $\beta_1$ en nous limitant à la classe des estimateurs linéaires.
Les estimateurs $\widehat{\widetilde{\beta}}_0$ et $\widehat{\widetilde{\beta}}_1$ sont linéaires s'ils sont de la forme 
\begin{align*}
    \widehat{\widetilde{\beta}}_0=\sum_{i=1}^n  \alpha_i Y_i, & & \widehat{\widetilde{\beta}}_1=\sum_{i=1}^n  \lambda_i Y_i,
\end{align*}
où $\alpha_1,\hdots,\alpha_n, \lambda_1,\hdots, \lambda_n$ sont des constantes connues.

\begin{parts}
    \part[4] Montrez que $\widehat{\widetilde{\beta}}_0$ est sans biais pour $\beta_0$ si les deux conditions suivantes sont respectées:
    \begin{align*}
        \sum_{i=1}^n  \alpha_i =1, &   & \sum_{i=1}^n  \alpha_i x_i =0.
    \end{align*}
    \part[4] Quelles conditions pouvons-nous imposer aux $\lambda_i$ pour que $\widehat{\widetilde{\beta}}_1$ soit sans biais pour $\beta_1$?
    \part[4] Montrez que l'estimateur linéaire non-biaisé de $\beta_1$ ayant la plus petite variance parmi tous les estimateurs linéaires non-biaisés de $\beta_1$ est celui qui minimise la valeur de $\displaystyle\sum_{i=1}^n\lambda_i ^2$.
    \part[2] Donnez l'expression de $ \alpha_i$, dans le cas de l'estimateur des moindres carrés pour $\beta_0$.
\end{parts}


\begin{solution}
    \begin{parts}
        \part $$\mathrm{E} = \sum \alpha_i (\beta_0 + \beta_1 x_i) = \beta_0 \sum \alpha_i + \beta_1 \sum \alpha_i x_i$$
        $ \sum \alpha_i = 1$ et $\sum \alpha_i x_i = 0$
        \part $$\mathrm{E} = \sum \lambda_i (\beta_0 + \beta_1 x_i)
        = \beta_0 \sum \lambda_i + \beta_1 \sum \lambda_i x_i$$
         $ \sum \lambda_i = 0$ et $\sum \lambda_i x_i = 1$
        \part $$ \mathrm{Var}  = \sum \lambda_i^2 \sigma^2$$
        \part $\widehat{\beta}_0 = \overline{Y} - \widehat{\beta}_1 \overline{x}$, les $d_i$ du cours.
    \end{parts}
\end{solution}